\section{Modelisation floue}

La modelisation floue V1 repose sur quatre variables linguistiques et une base minimale de huit regles. Les definitions sont externalisees dans \texttt{config/fuzzy\_config.yaml}, ce qui rend le systeme inspectable et ajustable sans modifier le code Python.

\subsection{Variables linguistiques V1}

\begin{table}[H]
\centering
\caption{Variables linguistiques de la V1.}
\label{tab:variables}
\small
\begin{tabular}{@{}lll@{}}
\toprule
Variable & Univers & Termes \\
\midrule
\texttt{genre\_preference} & $[0,1]$ & faible, moyenne, forte \\
\texttt{average\_rating} & $[0.5,5]$ & mauvaise, correcte, bonne, excellente \\
\texttt{popularity} & $[0,350]$ & confidentiel, modere, populaire, tres populaire \\
\texttt{recommendation\_score} & $[0,1]$ & tres faible, faible, moyen, fort, tres fort \\
\bottomrule
\end{tabular}
\end{table}

Les trois premieres variables sont les entrees du systeme Mamdani. La quatrieme est la variable de sortie defuzzifiee pour obtenir le score final.

\subsection{Fonctions d'appartenance}

\begin{figure}[H]
\centering
\includegraphics[width=0.82\textwidth]{membership_genre_preference.png}
\caption{Fonctions d'appartenance de la preference de genre.}
\label{fig:membership_genre_preference}
\end{figure}

\begin{figure}[H]
\centering
\includegraphics[width=0.82\textwidth]{membership_average_rating.png}
\caption{Fonctions d'appartenance de la note moyenne.}
\label{fig:membership_average_rating}
\end{figure}

\begin{figure}[H]
\centering
\includegraphics[width=0.82\textwidth]{membership_popularity.png}
\caption{Fonctions d'appartenance de la popularite.}
\label{fig:membership_popularity}
\end{figure}

\begin{figure}[H]
\centering
\includegraphics[width=0.82\textwidth]{membership_recommendation_score.png}
\caption{Fonctions d'appartenance du score de recommandation.}
\label{fig:membership_recommendation_score}
\end{figure}
\FloatBarrier

\subsection{Base de regles minimale}

\begin{table}[H]
\centering
\caption{Base de regles Mamdani minimale.}
\label{tab:rules}
\scriptsize
\begin{tabular}{@{}p{0.08\textwidth}p{0.59\textwidth}p{0.18\textwidth}@{}}
\toprule
Regle & Antecedents & Consequence \\
\midrule
R1 & preference forte, note excellente, popularite tres populaire & tres fort \\
R2 & preference forte, note excellente, popularite confidentiel & fort \\
R3 & preference forte, note bonne, popularite populaire & fort \\
R4 & preference moyenne, note excellente, popularite tres populaire & fort \\
R5 & preference moyenne, note bonne, popularite modere & moyen \\
R6 & preference faible, note excellente, popularite tres populaire & moyen \\
R7 & preference forte, note mauvaise, popularite populaire & faible \\
R8 & preference faible, note mauvaise, popularite confidentiel & tres faible \\
\bottomrule
\end{tabular}
\end{table}

Cette base de regles est volontairement compacte. Elle suffit a demontrer le raisonnement flou tout en gardant chaque decision explicable.

\subsection{Agregation et defuzzification}

Le moteur Mamdani utilise le minimum pour combiner les antecedents d'une regle, puis le maximum pour agreger les sorties activees. La defuzzification est realisee par centroide, conformement au parametrage \texttt{defuzzification\_method: centroid}. La figure d'inference de la section 9 illustre cette surface agregee et la position du centroide.

La section suivante precise comment les preferences utilisateur, souvent imprecises, sont projetees vers ces variables linguistiques.
